\chapter*{General Introduction}\label{chap:gen_intro}
\addcontentsline{toc}{chapter}{General Introduction}

Novel antimalarial therapies are required to combat drug-resistant \emph{Plasmodium falciparum} (Pf) and help with malaria elimination. Malaria is one of the deadliest disease afflicting the mankind, with more than $263$ million new cases every year, and over \num{597000} reported deaths (\emph{World Malaria Report} WHO\footnote{World Health Organization}, 2024). The disease poses a risk to half of the world's population, with non-Mediterranean Africa and Southeast Asia being particularly affected \cite{Alghamdi2024}. Malaria continues to represent a significant public health challenge across the African continent, particularly in regions situated south of the Sahara Desert, which accounted for $94\%$ of malaria cases and $95\%$ of malaria deaths in 2023. Notable reductions in incidence and mortality were observed between 2000 and 2022; however, the disease continues to disproportionately affect children under five, who represent a significant malaria fatalities. Key countries, including Nigeria, the Democratic Republic of Congo, Tanzania, and Niger, contribute substantially to the global mortality burden. In Cameroon, the entire population is susceptible to infection, with the highest risk concentrated in the equatorial and coastal regions. As one of the 11 countries with the highest malaria burden, Cameroon accounts for $2.6\%$ of all global malaria cases and $2.1\%$ of malaria deaths. The transmission of malaria is particularly pronounced in forest, coastal, and humid savanna regions, with the rainy season (July to October) contributing to an escalation in risk due to increased rainfall and temperature. The nation faces a multitude of interconnected challenges, including drug and insecticide resistance, climate change, and disparities in healthcare access, which collectively impede efforts to control malaria transmission by infected Anopheles mosquitoes. Moreover, malaria imposes a substantial economic burden, manifesting in increased healthcare expenditures, diminished productivity, and stunted GDP\footnote{Gross Domestic Product (GDP) is defined as the total market value of all goods and services produced within a nation's borders over a specific period. The fluctuations in GDP over time are typically measured on an annual or quarterly basis.} growth. This underscores the imperative for ongoing public health interventions and the development of novel control strategies.

In essence, a drug is any substance that has the capacity to cure or prevent diseases in humans and animals. For many centuries, drugs were derived from natural sources \cite{Schedler2006}. These substances were often discovered by chance, and their mechanisms of action were not always understood. Moreover, they were not always capable of being synthesized industrially. The process of discovering drugs was largely based on serendipity, or, in other words, fortunate accidents. The seminal contributions of Pasteur, who proposed the concept of a drug as a molecule with a biological mechanism of action, were informed by the preceding observations of van Leeuwenhoek. These observations, which included the identification of microorganisms and their association with diseases, laid the foundation for the field of microbiology. Indeed, at the beginning of the 20th century, Ehrlich articulated in \emph{"Principle of Chemotherapy"} that an appropriate dose of a chemical could disrupt the growth of disease-causing microorganisms \cite{Alizon2018}. This "principle" became a reality in the 1930s when Fleming discovered penicillin. Subsequent to this breakthrough, considerable efforts were made to synthesize or isolate compounds analogous to penicillin, which resulted in the development of various antibiotic families. Beginning in the 1950s, significant progress in the field of human genome research, complemented by substantial advancements in cell and molecular biology, paved the way for the development of rational drug design methodologies. In this novel approach, biologists formulate informed hypotheses concerning the mechanisms underlying diseases. Consequently, the process of drug discovery underwent a shift from a primarily chance-based approach to a more systematic one. This transformation occurred as researchers began to link diseases with specific proteins present within biological pathways, thereby facilitating the development of targeted organic compounds for therapeutic interventions. The objective of rational drug discovery is the identification of disease-associated targets and their ligands \cite{Kubinyi2003}. Consequently, in response to the AIDS epidemic that emerged in the 1980s, extensive research efforts culminated in the development of combinatorial chemistry, a field that facilitates the synthesis of a vast and diverse array of organic compounds. These compounds can then be tested for their efficacy in combating specific diseases. In the contemporary era, although chance still plays a role in the process of developing new drugs, recent technological advancements and discoveries have rendered the drug development process more rationalized \cite{Galloway2010}.

The process of drug discovery can be defined as the identification of novel pharmaceutical compounds with the potential to treat diseases \cite{Paul2010,Hughes2011}. In essence, this process entails the identification of novel chemical entities that possess specific chemical properties for the treatment of diseases. This process is often characterized by a significant investment of resources, both financial and temporal, and involves a rigorous testing and experimentation phase aimed at identifying drugs that are both safe and effective \cite{Schaduangrat2020,Sliwoski2014}. The primary objective of a scientist in the realm of drug discovery is the identification of lead or hit molecules of the highest quality. However, it is imperative that these lead molecules not only exhibit pharmacological activity but also ensure safety for both humans and the environment. Subsequent to the identification of a promising lead molecule, a series of trials and assessments are conducted to predict its pharmacokinetic properties, including Absorption, Distribution, Metabolism, Excretion, and Toxicity (ADMET) characteristics \cite{Schaduangrat2020,Lin2003,Naithani2024}. The process of developing new drugs is costly and time-consuming, necessitating the testing of thousands of chemical compounds and conducting numerous experiments to identify drugs that are both safe and effective. In light of these challenges, there is an imperative to establish robust, standard, and reproducible computational methodologies to achieve the objectives set. These computational approaches are instrumental in determining the characteristics of receptors and compounds, thereby facilitating the binding process \cite{Li2021}.

Despite progress in reducing malaria-related morbidity and mortality over the past two decades, the constant emergence of parasitic resistance to current antimalarial drugs has made the discovery of new drug candidates a major global health priority \cite{Neves2020,Raphemot2016}. Currently, artemisinin-combined therapy (ACT) since their implementation in 2006 remains the treatment of malaria, but Pf-resistance cases were recorded also. Malaria is known to have a large global burden, with strong poverty-promoting effects. It persists as chronic infections, in spite of the available effective medical treatments. While aggressive control measures have reduced prevalence rates in some regions, the emergence of drug resistance, particularly to artemisinin-based combination therapies, complicates treatment protocols and undermines elimination efforts. Therefore, the urgent need for innovative antimalarial compounds and therapeutic targets has become paramount \cite{Mathews2018}. In this context, understanding the sociocultural dynamics of health care partnerships in endemic regions can improve the effectiveness of interventions \cite{Kraeker2013}. Taken together, these challenges underscore the importance of sustained investment in research aimed at overcoming these barriers and achieving long-term malaria control.


Natural products (NPs) and their structural analogues have historically made significant contributions to pharmacotherapy, particularly for infectious diseases such as malaria. A notable example is  \emph{artemisinin}, a breakthrough treatment derived from \emph{Artemisia annua} (sweet wormwood), which was recognized with the Nobel Prize in Physiology or Medicine in 2015. However, natural products pose considerable challenges in the realm of drug discovery, primarily due to technical obstacles during the screening, isolation, characterization, and optimization processes. These challenges have contributed to a marked decline in the pharmaceutical industry's engagement in natural product research from the 1990s onwards \cite{Atanasov2021}. NPs exhibit distinctive characteristics when compared to conventional synthetic molecules, thus presenting both benefits and difficulties in the context of drug discovery for malaria. These compounds exhibit a remarkable diversity in their scaffolds and structural intricacy. NPs characteristically exhibit higher molecular weights, a greater number of sp3 carbon and oxygen atoms, and a lower number of nitrogen and halogen atoms in comparison to synthetic libraries. They also feature increased numbers of hydrogen bond acceptors and donors, demonstrate lower calculated octanol-water partition coefficients (indicating greater hydrophilicity), and exhibit enhanced molecular rigidity compared to synthetic alternatives. These distinctive properties can prove advantageous in the development of antimalarial drugs. For instance, the increased rigidity of NPs is valuable when targeting protein-protein interactions in Plasmodium parasites. Notably, natural products represent a significant source of oral antimalarial medications that function beyond Lipinski's rule of five \cite{Lipinski2004,Lipinski1997,Lipinski2004a}. The increasing molecular mass of approved oral medications, including antimalarials, over the last two decades demonstrates the growing importance of compounds that do not conform to this conventional pharmaceutical rule \cite{Shultz2018}.

Machine Learning (ML) have significantly advanced drug discovery. Pharmaceutical companies have greatly benefited from the utilization of various ML techniques to discover or re-purposing novel drugs candidate. The various ML algorithms have been used for predicting chemical, biological, and physical characteristics of compounds. Deep Learning (DL) techniques have become exemplar methods for drug activity prediction, target discovery and lead molecule discovery \cite{Patel2020}. These algorithms have been extremely used alongside the vast data generated utilizing high throughput sequencing to enhance the target discovery process. High throughput screening (HTS) using either approach entails screening of a large library of compounds. This process is often inefficient and not cost-effective because a high failure rate at subsequent stages of drug discovery. 

Biopharmaceutical industries are focusing on computational approaches in order to enhance the drug discovery processes as to reduce research and development expenses by diminishing failure rates in clinical trials and ultimately generate superior medicines. Drug discovery using ML techniques can be divided into four major steps: \emph{(i)} identify and curate datasets (potential target), \emph{(ii)} build and validate models, \emph{(iii)} score and prioritize a drug library, \emph{(iv)} validate predictions \emph{in vitro} or \emph{in vivo}. Different ML approaches help to identify drugs targets, find suitable molecules from data libraries, suggest chemical modifications, etc. The application of ML and Artificial Intelligence (AI) methods to drug discovery are reviewed in several recent papers. DL methods have made a massive impacts in science and technology generally, and in drug discovery particularly. Ashdown et \emph{al.} \cite{Ashdown2020} set out to develop a high-throughput cell-based, drug screening method that defines cell morphology and drug action from primary image data. From the outset, they sought to encompass asynchronous asexual cultures of the human malaria parasite Pf., by designing a semi-supervised DNN strategy that can analyze fluorescence microscopy images. The model can quantify the effects of antimalarial drugs both in terms of changes in life history distributions and morphological phenotype or drop points not observed during asexual development. Keshavarzi et \emph{al.} \cite{KeshavarziArshadi2020} introduce DL based process (\emph{DeepMalaria}), capable of predicting the anti-Pf inhibitory properties of compounds using a graph based model and their SMILES (Simplified Molecular Input Line Entry Specification) \cite{Weininger1988,Honda2019}. Monika Samant et \emph{al.} \cite{SamantMonika2016} have developed a protein-protein interaction network of Pf. and human host by integrating experimental data and computational prediction of interactions using the interlog method. To complement the medical professionals, Qazi Ammar et \emph{al.} \cite{Arshad2021} have devised a DL based method to automatically detect/localize the Pf. parasites in the photograph of stained film. They introduced a new large scale microscopic image malaria datasets, where cells are tagged from the microscopic images. Their model are able to classify Malaria life-cycle classification in thin blood smears images. Traditional machine learning approaches for drug discovery encounter computational limitations when processing complex molecular data, rendering quantum machine learning a potentially compelling alternative. This approach can potentially achieve exponential speedups, even with relatively modest quantum resources. The real question is, with all the modern technological advancements in drug discovery, how can we utilize innovative technologies to find a new active compounds more efficiently, thus reducing the cost?     

Quantum machine learning (QML) is an emerging field that integrates the principles of quantum computing with machine learning algorithms to potentially enhance computational capabilities. This innovative approach integrates the principles of quantum physics with the adaptability of machine learning, thereby offering new possibilities for solving complex problems across various domains \cite{Suzuki2020,Mensa2023}. The potential of the QML platform is significant; its integration into pharmaceutical development could result in substantial advancements, including enhanced molecular representation \cite{Roy2023,McKinseyCompany2023}, sophisticated feature extraction methods \cite{Huang2021}, accelerated screening processes \cite{Biamonte2017}, improved structure-activity relationship (SAR) analysis capabilities \cite{Cao2019}, optimized lead compound management \cite{Outeiral2020}, refined predictions of drug-target interactions \cite{Smaldone2024}, and streamlined handling of complex biological data \cite{PerdomoOrtiz2017,Cerezo2020}. The majority of QML applications in drug discovery are primarily focused on the initial stages of the drug discovery pipeline, particularly with regard to the identification of novel drug-like molecules. The complexity and intricacy of QML pose significant challenges in the interpretation of results, thereby hindering the identification of the molecular mechanisms that underpin drug-likeness \cite{Bova2021}. The utilization of QML models necessitates high-performance computing resources and specialized software \cite{Cerezo2023,Britt2015}. 

Quantum computing holds the potential to significantly enhance the drug discovery process by executing complex calculations with greater efficiency than classical computing systems. This capability is particularly advantageous during the initial phases of drug discovery, such as the identification of novel drug-like molecules and the prediction of molecular properties \cite{Batra2021,Avramouli2023,Avramouli2022,Batra2020}. By accelerating the research and development phase, QML can substantially reduce the costs associated with drug discovery. Some studies indicate that the integration of QML with quantum computing simulations could shorten the R\&D phase to mere months and decrease costs to a fraction of those incurred by traditional methods \cite{Wong2023}. QML frameworks, including those employing quantum support vector classifiers, have demonstrated superior performance in predicting the ADME-Tox properties of molecules compared to classical models \cite{Bhatia2023}. This enhancement is critical for informed decision-making within the pharmaceutical industry \cite{Bhatia2023}. A primary challenge in applying QML to drug discovery is the necessity of compressing extensive datasets of molecular descriptors to align with the capabilities of quantum computers, necessitating innovative data handling and algorithm design approaches \cite{Batra2021,Bhatia2023}. The amalgamation of quantum computing with machine learning is a complex endeavor requiring interdisciplinary collaboration. The integration process involves addressing technical challenges related to algorithm development and hardware limitations \cite{Avramouli2022}. Despite the potential benefits, QML remains in its early stages, and its practical application in drug discovery is constrained by the current state of quantum computing technology. Hybrid solutions that combine classical and quantum methodologies are presently recommended to capitalize on the strengths of both \cite{Avramouli2023}. QML techniques, such as generative adversarial networks and variational auto-encoders, are being investigated for the generation of small and large drug molecules and for virtual screening processes. These applications can improve the efficiency of identifying potential drug candidates \cite{Mensa2023,Li2021a}. The potential for QML to confer a quantum advantage in drug discovery is promising, particularly in ligand-based virtual screening and other computationally intensive tasks. Ongoing advancements in quantum hardware and algorithm development are anticipated to further augment these capabilities \cite{Mensa2023,Li2021a}.